\chapter*{Preface}
\addcontentsline{toc}{chapter}{Preface}
This work is a direct materialization of my hatred and rage, and was originally intended for not completely benign purposes. In short, it is a mathematical treatment of natural language. For a detailed description, keep reading.

The idea of approaching languages mathematically might seem odd, but has more merit than it lets in on. A language is more than merely a set of symbols: it is a perspective of perceiving the world. A real learner will never be satisfied by merely knowing how to speak or understand, but will realize the true meaning of what is said. This state of absolute mastery is only possible if one embodies the soul of the target language rather than merely translating phrases, which requires knowledge of its structural finesse: something  people normally call `grammar'. Hence, a comprehensive assessment of target grammar against source grammar will allow a seamless and simple transition into the new perspective.

Unfortunately, no such assessment exists. Worse, typical language courses as well as most people treat language as something non-technical, and leave the acquisition process entirely to the subconscious mind. This is similar to trying to learn a programming language only by looking at code, and without knowing its syntax: a waste of time. Also, most grammar explanations have to be integrated from multiple sources to be complete, a task whose inherent complexity defeats the purpose. This is not a problem I appreciate, and thus lies on my chopping block at the moment. If nothing else, it is something I can joyously murder for the rest of my life without reaching stagnation. The motivation explained here is considered extensively in \autoref{chap:int}, which also contains important points about the structure of this work, and \hyperref[copyright]{\textcolor{magenta}{information on copyright}}.

Other reasons for creating this series: I wish to allow people to bridge and integrate into other cultures without having to use English as a medium,  facilitate their understanding of said cultures, show that languages are in fact technical and mostly logical, prove that at their core; all languages are the exact same, help learners save unnecessary effort,  inspire other people to create their own structural transformations in the future, and because I believe that it is only through understanding that real hatred can be born.

\textit{This work has no original language, since it is meant to be a `solver'. However, since I have limitations on how well I may master all languages I seek to solve, only the German and English versions will have a detailed introduction. The solver is fluid, and readers are encouraged to keep updating it according to their experiences and requirements.}